% Source: Hatcher, "Algebraic Topology", Chapter 0 "Some Underlying Geometric Notions"
% (2002 Cambridge University Press edition), PDF pages 10-29 of the cited source.
% Transcribed from the cited source and organized according to
% leanblueprint conventions (semantic \label, bracketed titles, \uses dependency edges) below.

Maps between spaces are assumed continuous unless explicitly stated otherwise.

\section{Homotopy and Homotopy Type}

\begin{definition}[Deformation retraction]
\label{def:deformation-retraction}
\group{Homotopy Theory}
\lean{HatcherLib.DeformationRetract}
\leanok
\uses{def:retraction,def:homotopy-rel-a}
A \textbf{deformation retraction} of a space $X$ onto a subspace $A$ is a family
$f_t:X\to X$ for $t\in I=[0,1]$ such that $f_0=\one$, $f_1(X)=A$, and
$f_t|_A=\one$ for every $t$, with $(x,t)\mapsto f_t(x)$ continuous.
\end{definition}

\begin{definition}[Mapping cylinder]
\label{def:mapping-cylinder}
\group{Homotopy Theory}
\lean{HatcherLib.MappingCylinder}
\leanok
For a map $f:X\to Y$, the \textbf{mapping cylinder} $M_f$ is the quotient of
$(X\times I)\sqcup Y$ by $(x,1)\sim f(x)$.
\end{definition}

The mapping cylinder deformation retracts onto $Y$ by moving $[x,t]$ to $[x,1]=f(x)$.

\begin{definition}[Homotopy]
\label{def:homotopy}
\group{Homotopy Theory}
\lean{HatcherLib.Homotopy}
\leanok
A \textbf{homotopy} from maps $f_0,f_1:X\to Y$ is a family $f_t:X\to Y$ whose
adjoint $F:X\times I\to Y$, $F(x,t)=f_t(x)$, is continuous.
\end{definition}

\begin{definition}[Homotopic maps]
\label{def:homotopic-maps}
\group{Homotopy Theory}
\lean{HatcherLib.Homotopic}
\leanok
\uses{def:homotopy}
Maps $f_0,f_1:X\to Y$ are \textbf{homotopic}, written $f_0\simeq f_1$, if a homotopy
connects them.
\end{definition}

\begin{definition}[Retraction]
\label{def:retraction}
\group{Homotopy Theory}
\lean{HatcherLib.IsRetraction}
\leanok
A \textbf{retraction} of $X$ onto $A\subset X$ is a map $r:X\to X$ with
$r(X)=A$ and $r|_A=\one$. Equivalently, $r^2=r$; one may instead regard $r$ as
a map $X\to A$ restricting to the identity on $A$.
\end{definition}

A deformation retraction is a homotopy rel $A$ from $\one_X$ to a retraction.

\begin{definition}[Homotopy relative to a subspace]
\label{def:homotopy-rel-a}
\group{Homotopy Theory}
\lean{HatcherLib.HomotopyRel}
\leanok
\uses{def:homotopy}
A homotopy $f_t:X\to Y$ is \textbf{relative to $A$}, or \textbf{rel $A$}, if
$f_t|_A$ is independent of $t$.
\end{definition}

If $X$ deformation retracts onto $A$ via $f_t$, its terminal map factors through a
retraction $r:X\to A$ and the inclusion $i:A\hookrightarrow X$ satisfies
$ri=\one_A$ and $ir\simeq\one_X$.

\begin{definition}[Homotopy equivalence and homotopy type]
\label{def:homotopy-equivalence}
\group{Homotopy Theory}
\lean{HatcherLib.HomotopyEquiv}
\leanok
\uses{def:homotopy,def:homotopic-maps}
A map $f:X\to Y$ is a \textbf{homotopy equivalence} if some $g:Y\to X$ satisfies
$fg\simeq\one_Y$ and $gf\simeq\one_X$. In this case $X$ and $Y$ have the same
\textbf{homotopy type}, written $X\simeq Y$.
\end{definition}

Homotopy equivalence is an equivalence relation on spaces. Composites of homotopy
equivalences are homotopy equivalences, and a map homotopic to a homotopy
equivalence is itself a homotopy equivalence.

\begin{proposition}[A deformation retraction is a homotopy equivalence]
\label{prop:deformation-retract-homotopy-equivalence}
\group{Homotopy Theory}
\lean{HatcherLib.DeformationRetract.homotopyEquiv}
\leanok
\uses{def:deformation-retraction,def:homotopy-equivalence}
If $X$ deformation retracts onto $A$, then $A\simeq X$ and the inclusion
$i:A\hookrightarrow X$ is a homotopy equivalence.
\end{proposition}

\begin{proof}
\leanok
\uses{def:retraction}
\sketch Let $r=f_1:X\to X$ be the terminal retraction and corestrict it to
$\bar r:X\to A$. Then $\bar r i=\one_A$, while $i\bar r=r=f_1\simeq f_0=\one_X$.
Thus $i$ and $\bar r$ are homotopy inverses.
\end{proof}

\begin{definition}[Contractible space]
\label{def:contractible}
\group{Homotopy Theory}
\lean{HatcherLib.IsContractible}
\leanok
\uses{def:homotopy-equivalence}
A space is \textbf{contractible} if it has the homotopy type of a point, equivalently
if its identity map is nullhomotopic. This is weaker in general than deformation
retracting to a point.
\end{definition}

\begin{example}[The house with two rooms is contractible]
\label{ex:house-with-two-rooms-contractible}
\group{Homotopy Theory}
\uses{def:contractible,def:mapping-cylinder}
The \textit{house with two rooms} is the 2-dimensional subspace $X\subset\R^3$
formed from three horizontal annuli and vertical rectangular chamber and tunnel walls.
For sufficiently small $\varepsilon$, a closed neighborhood $N(X)$ deformation retracts
onto $X$ and is the mapping cylinder of its boundary surface mapping to $X$.
\end{example}

\begin{proof}
\sketch A family of embeddings of $D^3$ into $\R^3$, obtained by creating the two
tunnels and chambers, identifies $N(X)$ with $D^3$. Hence
$X\simeq N(X)\cong D^3\simeq\text{point}$. If $f_t$ contracts $N(X)$ to
$x_0\in X$ and $r:N(X)\to X$ is a retraction, then $(r\circ f_t)|_X$ contracts
$X$ to $x_0$.
\end{proof}

\section{Cell Complexes}
\label{sec:cell-complexes}

\begin{definition}[CW complex]
\label{def:cell-complex}
\group{CW Complexes}
\lean{HatcherLib.IsCWComplex}
\leanok
A \textbf{cell complex} or \textbf{CW complex} is built inductively as follows:
\begin{enumerate}
\item[(1)] $X^0$ is a discrete set of 0-cells.
\item[(2)] $X^n$ is obtained from $X^{n-1}$ by attaching disks $D^n_\alpha$ along
maps $\varphi_\alpha:S^{n-1}\to X^{n-1}$. Thus
$X^n=X^{n-1}\sqcup_\alpha e^n_\alpha$ as a set, where $e^n_\alpha$ is an open disk.
\item[(3)] The construction either stops at a finite stage $X=X^n$, or
$X=\bigcup_nX^n$ with the weak topology: $B\subset X$ is open or closed exactly
when $B\cap X^n$ has that property in every $X^n$.
\end{enumerate}
\end{definition}

\begin{definition}[Dimension of a CW complex]
\label{def:dimension-cw-complex}
\group{CW Complexes}
\uses{def:cell-complex}
If $X=X^n$ and $n$ is minimal, then $X$ is finite-dimensional of dimension $n$.
\end{definition}

\begin{example}[Graphs as 1-dimensional CW complexes]
\label{ex:cw-graph}
\dcref{ch0:0.1}
\group{CW Complexes}
\mathlibok
\uses{def:cell-complex}
A 1-dimensional CW complex is a graph: its 0-cells are vertices and its 1-cells
are edges, whose two ends may attach to the same vertex.
\end{example}

\begin{example}[Euler characteristic of the house with two rooms]
\label{ex:house-two-rooms-euler-characteristic}
\dcref{ch0:0.2}
\group{CW Complexes}
\uses{ex:house-with-two-rooms-contractible,def:cell-complex,thm:euler-characteristic-homology}
The house with two rooms has a 2-dimensional CW structure with 29 0-cells, 51 1-cells,
and 23 2-cells. Its Euler characteristic is
$\chi(X)=29-51+23=1$, where for a finite CW complex $\chi$ is the number of
even-dimensional cells minus the number of odd-dimensional cells. Its contractibility
therefore agrees with homotopy invariance of Euler characteristic.
\end{example}

\begin{example}[Minimal CW structure on the sphere $S^n$]
\label{ex:sphere-two-cells}
\dcref{ch0:0.3}
\group{CW Complexes}
\uses{def:cell-complex}
$S^n$ has two cells, $e^0$ and $e^n$, with the $n$-cell attached by the constant
map $S^{n-1}\to e^0$; equivalently, $S^n=D^n/\partial D^n$.
\end{example}

\begin{example}[CW structure on real projective space $\RP^n$]
\label{ex:real-projective-space-cw}
\dcref{ch0:0.4}
\group{CW Complexes}
\uses{def:cell-complex,ex:sphere-two-cells}
$\RP^n$ is the quotient of $S^n$ by $v\sim -v$, equivalently a hemisphere $D^n$
with antipodal points of its boundary identified. Thus $\RP^n$ is obtained from
$\RP^{n-1}$ by attaching an $n$-cell via $S^{n-1}\to\RP^{n-1}$, and has one cell
$e^i$ in each dimension $0\le i\le n$.
\end{example}

\begin{example}[CW structure on $\RP^\infty$]
\label{ex:infinite-real-projective-space}
\dcref{ch0:0.5}
\group{CW Complexes}
\uses{ex:real-projective-space-cw}
$\RP^\infty=\bigcup_n\RP^n$ is a CW complex with one cell in every dimension;
it is the space of lines in $\R^\infty=\bigcup_n\R^n$.
\end{example}

\begin{example}[CW structure on complex projective space $\CP^n$]
\label{ex:complex-projective-space-cw}
\dcref{ch0:0.6}
\group{CW Complexes}
\uses{def:cell-complex}
$\CP^n$ is the space of complex lines in $\C^{n+1}$. Viewing it as the quotient of
the unit sphere in $\C^{n+1}$ by multiplication by unit complex scalars, or as a disk
$D^{2n}$ with its boundary identified under this action, shows that $\CP^n$ is obtained
from $\CP^{n-1}$ by attaching a $2n$-cell. Hence
$\CP^n=e^0\cup e^2\cup\cdots\cup e^{2n}$, and $\CP^\infty$ has one cell in each
even dimension.
\end{example}

\begin{definition}[Characteristic map of a cell]
\label{def:characteristic-map}
\group{CW Complexes}
\uses{def:cell-complex}
Each cell $e^n_\alpha$ has a \textbf{characteristic map}
$\Phi_\alpha:D^n_\alpha\to X$ extending its attaching map and restricting to a
homeomorphism from the disk interior onto $e^n_\alpha$.
\end{definition}

\begin{definition}[Subcomplex]
\label{def:subcomplex}
\group{CW Complexes}
\mathlibok
\uses{def:cell-complex}
A \textbf{subcomplex} is a closed subspace $A\subset X$ that is a union of cells.
The attaching image of every cell in $A$ lies in $A$, so $A$ is itself a CW complex.
\end{definition}

\begin{definition}[CW pair]
\label{def:cw-pair}
\group{CW Complexes}
\mathlibok
\uses{def:subcomplex}
A pair $(X,A)$ consisting of a CW complex and a subcomplex is a \textbf{CW pair}.
\end{definition}

The closure of a cell need not be a subcomplex: attach a 2-cell to the minimal CW
structure on $S^1$ by a map whose image is a proper nontrivial arc; the 2-cell closure
contains only part of the 1-cell.

\begin{remark}[Closure of a cell need not be a subcomplex]
\label{rem:closure-of-cell-not-always-subcomplex}
\group{CW Complexes}
\uses{def:subcomplex}
The preceding construction shows that a cell closure, and more generally the closure
of a union of cells, need not be a subcomplex. In many standard CW structures,
however, the closure of every union of cells is a subcomplex.
\end{remark}

\section{Operations on Spaces}
\label{sec:operations-on-spaces}

\begin{definition}[Product CW structure]
\label{def:product-cw-structure}
\group{CW Complexes}
\uses{def:cell-complex}
For CW complexes $X$ and $Y$, the product CW structure has cells
$e^m_\alpha\times e^n_\beta$. Its CW topology can be finer than the product topology;
the topologies agree if either factor has finitely many cells or both have countably
many cells.
\end{definition}

\begin{definition}[Quotient CW structure]
\label{def:quotient-cw-structure}
\group{CW Complexes}
\uses{def:cw-pair}
For a CW pair $(X,A)$, $X/A$ has the cells of $X-A$ together with one new 0-cell.
The attaching map of a remaining cell is the composite of its original attaching map
with $X^{n-1}\to X^{n-1}/A^{n-1}$.
\end{definition}

\begin{definition}[Suspension]
\label{def:suspension}
\group{Homotopy Theory}
For a space $X$, its \textbf{suspension} $SX$ is
$(X\times I)/(X\times\{0\})/(X\times\{1\})$, with the two collapsed sets becoming
distinct suspension points. In particular $SS^n\cong S^{n+1}$, and a map $f:X\to Y$
induces $Sf:SX\to SY$ from $f\times I$. If $X$ is a CW complex, so are $SX$ and
the cone below.
\end{definition}

\begin{definition}[Cone]
\label{def:cone}
\group{Homotopy Theory}
The \textbf{cone} on $X$ is $CX=(X\times I)/(X\times\{0\})$; $SX$ is the union
of two copies of $CX$.
\end{definition}

\begin{definition}[Join of spaces]
\label{def:join}
\group{Homotopy Theory}
\uses{def:cone,def:suspension}
The \textbf{join} $X*Y$ is the quotient of $X\times Y\times I$ by
$(x,y_1,0)\sim(x,y_2,0)$ and $(x_1,y,1)\sim(x_2,y,1)$. Its points can be written
as formal combinations $t_1x+t_2y$ with $t_i\ge0$, $t_1+t_2=1$, subject to
$0x+1y=y$ and $1x+0y=x$.
\end{definition}

\begin{definition}[Simplex]
\label{def:simplex}
\group{Homotopy Theory}
\uses{def:join}
The join of $n$ points is the $(n-1)$-simplex
\[
\Delta^{n-1}=\{(t_1,\ldots,t_n)\in\R^n\mid t_i\ge0,\ \sum_i t_i=1\}.
\]
The join of $n$ copies of $S^0$ is homeomorphic, by radial projection, to
$S^{n-1}$.
\end{definition}

If $X$ and $Y$ are CW complexes, $X*Y$ has a natural CW structure with $X$ and $Y$
as subcomplexes and remaining cells $e^m_\alpha\times e^n_\beta\times(0,1)$.
Its CW topology can be finer than the quotient of the product topology on
$X\times Y\times I$.

\begin{definition}[Wedge sum]
\label{def:wedge-sum}
\group{Homotopy Theory}
For pointed spaces $(X,x_0)$ and $(Y,y_0)$, the \textbf{wedge sum} $X\vee Y$ is
$(X\sqcup Y)/(x_0\sim y_0)$. An arbitrary wedge of CW complexes is a CW complex
when every chosen basepoint is a $0$-cell. Moreover, $X^n/X^{n-1}$ is the wedge of
one $S^n$ for each $n$-cell.
\end{definition}

\begin{definition}[Smash product]
\label{def:smash-product}
\group{Homotopy Theory}
\uses{def:wedge-sum}
For pointed spaces, the \textbf{smash product} is
$X\wedge Y=(X\times Y)/(X\vee Y)$, where
$X\vee Y=(X\times\{y_0\})\cup(\{x_0\}\times Y)$. If $X$ and $Y$ are CW
complexes and $x_0,y_0$ are $0$-cells, this is a CW complex when equipped with
the CW topology, which must be used instead of the quotient/product topology
when the two differ. Also $S^m\wedge S^n\cong S^{m+n}$.
\end{definition}

\section{Two Criteria for Homotopy Equivalence}
\label{sec:two-criteria-homotopy-equivalence}

\begin{proposition}[Collapsing a contractible CW subcomplex]
\label{prop:collapsing-contractible-cw-subcomplex}
\group{CW Complexes}
\uses{def:cw-pair,def:contractible}
If $(X,A)$ is a CW pair and $A$ is contractible, then the quotient map
$X\to X/A$ is a homotopy equivalence.
\end{proposition}

\begin{definition}[Attaching one space to another]
\label{def:attaching-space}
\group{CW Complexes}
\lean{HatcherLib.AttachingSpace}
\leanok
Given $A\subset X_1$, $X_0\sqcup_fX_1$ denotes the quotient of $X_0\sqcup X_1$
by $a\sim f(a)$ for a map $f:A\to X_0$. For $(X_1,A)=(D^n,S^{n-1})$ this is
attachment of an $n$-cell.
\end{definition}

The mapping cylinder is obtained by attaching $X\times I$ to $Y$ along
$X\times\{1\}$, and the mapping cone is the analogous construction with the cone.

\begin{definition}[Mapping cone]
\label{def:mapping-cone}
\group{Homotopy Theory}
\uses{def:cone,def:attaching-space}
For $f:X\to Y$, the \textbf{mapping cone} is
$C_f=Y\sqcup_f CX$, attaching $X\times\{1\}$ to $Y$ by $f$. If
$X=S^{n-1}$, this attaches an $n$-cell to $Y$; also $C_f=M_f/X$.
\end{definition}

\begin{proposition}[Homotopic attaching maps yield homotopy equivalent spaces]
\label{prop:attaching-homotopic-maps-preview}
\group{CW Complexes}
\uses{def:cw-pair,def:attaching-space}
If $(X_1,A)$ is a CW pair and $f,g:A\to X_0$ are homotopic, then
$X_0\sqcup_fX_1\simeq X_0\sqcup_gX_1$.
\end{proposition}

\begin{remark}[Sharper version of homotopic attaching maps]
\label{rem:attaching-homotopic-maps-sharper-version}
\group{CW Complexes}
\uses{prop:attaching-maps-homotopic-rel-x0}
The later proposition \cref{prop:attaching-maps-homotopic-rel-x0} gives a relative
$X_0$ version with a proof.
\end{remark}

\begin{example}[Quotient by a subcomplex as a mapping cone]
\label{ex:quotient-contractible-subcomplex-mapping-cone}
\dcref{ch0:0.13}
\group{Homotopy Theory}
\uses{def:cw-pair,def:mapping-cone}
For a CW pair $(X,A)$, the quotient $X/A$ is homotopy equivalent to the mapping cone
$X\cup CA$ of $A\hookrightarrow X$, since
$X/A=(X\cup CA)/CA$ and $CA$ is contractible.
\end{example}

\begin{example}[Quotient by a subspace contractible in $X$]
\label{ex:quotient-by-contractible-in-x}
\dcref{ch0:0.14}
\group{Homotopy Theory}
\uses{ex:quotient-contractible-subcomplex-mapping-cone,def:suspension}
If $(X,A)$ is a CW pair and $A\hookrightarrow X$ is homotopic to a constant map, then
$X/A\simeq X\vee SA$. In particular, $S^n/S^i\simeq S^n\vee S^{i+1}$ for $i<n$.
\end{example}

\begin{example}[Reduced suspension]
\label{ex:reduced-suspension}
\dcref{ch0:0.10}
\group{Homotopy Theory}
\uses{def:suspension,def:smash-product}
If $X$ is a CW complex with 0-cell $x_0$, collapsing the segment
$\{x_0\}\times I\subset SX$ gives the reduced suspension $\Sigma X\simeq SX$.
For CW complexes $X$ and $Y$ with chosen $0$-cells, it satisfies
$\Sigma(X\vee Y)=\Sigma X\vee\Sigma Y$ and $\Sigma X=X\wedge S^1$.
For example, $\Sigma(S^1\vee S^1)=S^2\vee S^2$.
\end{example}

\section{The Homotopy Extension Property}
\label{sec:homotopy-extension-property}

\begin{definition}[Homotopy extension property]
\label{def:homotopy-extension-property}
\group{CW Complexes}
\lean{HatcherLib.HasHEP}
\leanok
\uses{def:homotopy}
The pair $(X,A)$ has the \textbf{homotopy extension property} if every pair of maps
$X\times\{0\}\to Y$ and $A\times I\to Y$ that agree on $A\times\{0\}$ extends
to a map $X\times I\to Y$.
\end{definition}

\begin{proposition}[HEP is equivalent to a retraction property]
\label{prop:hep-iff-retract}
\group{CW Complexes}
\lean{HatcherLib.hasHEP_iff_isRetract}
\leanok
\uses{def:homotopy-extension-property,def:retraction}
If $A$ is closed in $X$, then $(X,A)$ has HEP iff
$X\times\{0\}\cup A\times I$ is a retract of $X\times I$. The forward implication
does not require closedness; the converse for arbitrary $A$ requires the appendix
argument. If $X$ is Hausdorff and such a retraction exists, then $A$ is closed.
\end{proposition}

\begin{proof}
\leanok
\sketch HEP extends the identity of
$X\times\{0\}\cup A\times I$. Conversely, when $A$ is closed, compatible maps on
the two closed pieces paste to a map on their union, which composes with the retraction
to give the required extension. In a Hausdorff $X$, the image of a retraction is its
closed fixed-point set, so the cylinder subspace is closed and hence $A$ is closed.
\end{proof}

\begin{lemma}[A retract $X \times \{0\} \cup A \times I$ of the cylinder is a deformation retract]
\label{lem:hep-cylinder-deformation-retract}
\group{CW Complexes}
\lean{HatcherLib.hepBaseDeformationRetract, HatcherLib.HasHEP.hepBase_deformationRetract}
\leanok
\uses{def:homotopy-extension-property,def:deformation-retraction,prop:hep-iff-retract}
If $r:X\times I\to X\times\{0\}\cup A\times I$ is a retraction, then the target
is a deformation retract. Consequently HEP implies this deformation retraction.
\end{lemma}

\begin{proof}
\leanok
\sketch Write $r(x,s)=(r_1(x,s),r_2(x,s))$ and set
\[
h_u(x,s)=\bigl(r_1(x,us),(1-u)s+u r_2(x,s)\bigr).
\]
This fixes $X\times\{0\}\cup A\times I$, starts at the identity, and ends at $r$.
\end{proof}

\begin{example}[A pair without the homotopy extension property]
\label{ex:hep-fails-for-bad-subspace}
\group{CW Complexes}
\uses{prop:hep-iff-retract}
Let $A=\{0,1,1/2,1/4,1/3,\ldots\}\subset I$. Although $A$ is closed, $(I,A)$
does not have HEP: there is no retraction
$I\times I\to I\times\{0\}\cup A\times I$.
\end{example}

\begin{example}[Mapping cylinder neighborhoods give the HEP]
\label{ex:mapping-cylinder-neighborhood-hep}
\dcref{ch0:0.15}
\group{CW Complexes}
\uses{def:mapping-cylinder,prop:hep-iff-retract}
A pair $(X,A)$ has HEP if $A$ has a mapping cylinder neighborhood: a closed
neighborhood $N$ with boundary part $B$, an open neighborhood $N-B$ of $A$, and a
homeomorphism $M_f\cong N$ fixed on $A\cup B$ for some $f:B\to A$.
\end{example}

\begin{proof}
\sketch The standard retraction of $I\times I$ onto
$I\times\{0\}\cup\partial I\times I$ gives HEP for $(M_f,A\cup B)$ and hence for
$(N,A\cup B)$. Extend a prescribed homotopy over $N$ and use the constant homotopy
outside $N-B$.
\end{proof}

\begin{proposition}[CW pairs have the homotopy extension property]
\label{prop:cw-pairs-have-hep}
\dcref{ch0:0.16}
\group{CW Complexes}
\uses{def:cw-pair,def:homotopy-extension-property,def:deformation-retraction}
If $(X,A)$ is a CW pair, then
$X\times\{0\}\cup A\times I$ is a deformation retract of $X\times I$; hence
$(X,A)$ has HEP.
\end{proposition}

\begin{proof}
\sketch Radial projection gives a deformation retraction
$D^n\times I\to D^n\times\{0\}\cup\partial D^n\times I$. Apply this on each
$n$-cell, extending over the skeleta. Perform the $n$-stage during
$[2^{-(n+1)},2^{-n}]$; the resulting infinite concatenation is continuous by the
weak topology and gives the stated deformation retraction.
\end{proof}

\begin{proposition}[Quotient by a contractible subspace with the HEP]
\label{prop:quotient-by-contractible-hep-homotopy-equiv}
\dcref{ch0:0.17}
\group{CW Complexes}
\lean{HatcherLib.collapseMk_homotopyEquiv}
\leanok
\uses{def:homotopy-extension-property,def:contractible,def:homotopy-equivalence}
If $(X,A)$ has HEP and $A$ is contractible, then the quotient map $q:X\to X/A$ is
a homotopy equivalence.
\end{proposition}

\begin{proof}
\leanok
\sketch Extend a contraction of $A$ to $f_t:X\to X$ with $f_0=\one$. Since
$f_t(A)\subset A$, $qf_t$ descends to $\bar f_t:X/A\to X/A$. At $t=1$, $f_1$
is constant on $A$ and induces $g:X/A\to X$ with $gq=f_1$. Thus
$gq=f_1\simeq\one_X$ and $qg=\bar f_1\simeq\one_{X/A}$.
\end{proof}

\begin{lemma}[Homotopic attaching maps and the HEP]
\label{lem:attach-homotopic-hep-rel}
\group{CW Complexes}
\lean{HatcherLib.attachingSpace_homotopyEquiv_rel, HatcherLib.attachingSpace_homotopyEquiv}
\leanok
\uses{def:attaching-space,def:homotopy-extension-property,lem:hep-cylinder-deformation-retract,def:homotopy-rel-a}
Let $A\subset X_1$ be closed, let $(X_1,A)$ have HEP, and let homotopic maps
$f,g:A\to X_0$ be attaching maps. Then
$X_0\sqcup_fX_1\simeq X_0\sqcup_gX_1\ \rel X_0$.
\end{lemma}

\begin{proof}
\leanok
\sketch For a homotopy $F:A\times I\to X_0$, form
$W=X_0\sqcup_F(X_1\times I)$. HEP gives deformation retractions of $W$ onto
$X_0$ together with either endpoint level $X_1\times\{0\}$ or
$X_1\times\{1\}$, relative to $X_0$. Closedness of $A$ makes the collapse and
pasting maps defining these retractions continuous. Composing the endpoint
inclusions and retractions gives homotopy inverses between
$X_0\sqcup_fX_1$ and $X_0\sqcup_gX_1$ relative to $X_0$.
\end{proof}

\begin{proposition}[Homotopic attaching maps give a homotopy equivalence rel $X_0$]
\label{prop:attaching-maps-homotopic-rel-x0}
\dcref{ch0:0.18}
\group{CW Complexes}
\uses{def:cw-pair,def:attaching-space,prop:cw-pairs-have-hep,lem:attach-homotopic-hep-rel}
If $(X_1,A)$ is a CW pair and $f,g:A\to X_0$ are homotopic, then
$X_0\sqcup_fX_1\simeq X_0\sqcup_gX_1\ \rel X_0$.
\end{proposition}

\begin{proof}
\sketch A subcomplex is closed and the preceding proposition applies because CW
pairs have HEP.
\end{proof}

\begin{lemma}[The homotopy extension property is preserved by products with $I$]
\label{lem:hep-product-interval}
\group{CW Complexes}
\lean{HatcherLib.HasHEPMap.prodMap_id}
\leanok
\uses{def:homotopy-extension-property}
If $(X,A)$ has HEP, then $(X\times I,A\times I)$ has HEP.
\end{lemma}

\begin{proof}
\leanok
\sketch Use the exponential correspondence between maps
$X\times I\times I\to Y$ and maps $X\times I\to C(I,Y)$. The extension problem
for $(X\times I,A\times I)$ is the HEP extension problem for $(X,A)$ with target
$C(I,Y)$, and evaluation recovers the required extension.
\end{proof}

\begin{proposition}[A homotopy equivalence restricting to the identity on $A$ is a homotopy equivalence rel $A$]
\label{prop:hep-homotopy-equivalence-rel-a}
\dcref{ch0:0.19}
\group{CW Complexes}
\lean{HatcherLib.hep_homotopy_equiv_rel}
\leanok
\uses{def:homotopy-extension-property,def:homotopy-equivalence,lem:hep-product-interval}
Suppose $(X,A)$ and $(Y,A)$ have HEP and $f:X\to Y$ is a homotopy equivalence with
$f|_A=\one$. Then $f$ is a homotopy equivalence rel $A$.
\end{proposition}

\begin{proof}
\leanok
\sketch Let $g$ be a homotopy inverse. Extend the restriction to $A$ of a homotopy
$gf\simeq\one_X$ to change $g$ within its homotopy class to $g_1$ with
$g_1|_A=\one$. Concatenate the resulting homotopy $g_1f\simeq\one_X$ with the
original homotopy, and use HEP for $(X\times I,A\times I)$ to make it rel $A$.
Repeat the construction with $(g_1,f)$ to obtain $f_1|_A=\one$ and
$f_1g_1\simeq\one$ rel $A$. Then $f_1\simeq f$ rel $A$ and
$fg_1\simeq\one$ rel $A$.
\end{proof}

\begin{corollary}[HEP inclusion homotopy equivalence implies deformation retract]
\label{cor:hep-inclusion-homotopy-equiv-deformation-retract}
\dcref{ch0:0.20}
\group{CW Complexes}
\lean{HatcherLib.hep_inclusion_deformation_retract}
\leanok
\uses{prop:hep-homotopy-equivalence-rel-a}
If $(X,A)$ has HEP and $A\hookrightarrow X$ is a homotopy equivalence, then $A$ is
a deformation retract of $X$.
\end{corollary}

\begin{proof}
\leanok
\sketch Apply the proposition to the inclusion of pairs $(A,A)\to(X,A)$. Its
relative homotopy inverse is a retraction $r:X\to A$ with $\iota r\simeq\one_X$
rel $A$, which is a deformation retraction.
\end{proof}

\begin{lemma}[The pair $(M_f, X)$ has the homotopy extension property]
\label{lem:mapping-cylinder-pair-hep}
\group{CW Complexes}
\lean{HatcherLib.hasHEPMap_mcylInclX}
\leanok
\uses{def:mapping-cylinder,def:homotopy-extension-property}
For any $f:X\to Y$, the pair $(M_f,X)$, with $X=X\times\{0\}\subset M_f$, has HEP.
\end{lemma}

\begin{proof}
\leanok
\sketch For a map $\varphi:M_f\to Z$ and a homotopy $h:X\times I\to Z$ of its
restriction, project $I\times I$ from $(1,2)$ onto the bottom and left edges. The
projection is
\[
\alpha(s,t)=\frac{2s-t}{2-t}\quad(t\le2s),\qquad
\beta(s,t)=\frac{t-2s}{1-s}\quad(t\ge2s).
\]
Define $F([x,s],t)=\varphi[x,\alpha(s,t)]$ in the first region and
$F([x,s],t)=h(x,\beta(s,t))$ in the second, while setting $F([y],t)=\varphi[y]$.
The formulas agree on $t=2s$, respect the mapping-cylinder quotient, and satisfy
$F(\cdot,0)=\varphi$ and $F([x,0],t)=h(x,t)$.
\end{proof}

\begin{corollary}[Homotopy equivalences are exactly deformation retractions of mapping cylinders]
\label{cor:homotopy-equivalence-iff-deformation-retract-mapping-cylinder}
\dcref{ch0:0.21}
\group{CW Complexes}
\lean{HatcherLib.isHmtpyEquiv_iff_isMcylDeformationRetract, HatcherLib.homotopyEquiv_iff_common_deformationRetract}
\leanok
\uses{cor:hep-inclusion-homotopy-equiv-deformation-retract,def:mapping-cylinder,lem:mapping-cylinder-pair-hep}
A map $f:X\to Y$ is a homotopy equivalence iff $X$ is a deformation retract of
$M_f$. Consequently, $X\simeq Y$ iff some space contains both as deformation retracts.
\end{corollary}

\begin{proof}
\leanok
\sketch Let $i:X\to M_f$ be the bottom inclusion and $r:M_f\to Y$ the canonical
retraction, so $f=ri$. Since $r$ is a homotopy equivalence, $f$ is one iff $i$ is:
the converse uses the inclusion $j:Y\to M_f$ and $i\simeq jri=jf$. If $X$ is a
deformation retract, then $i$ and hence $f$ are equivalences. Conversely, if $f$ is
one, so is $i$, and the HEP of $(M_f,X)$ upgrades a homotopy inverse of $i$ to a
deformation retraction. This uses closure of homotopy equivalences under composition
and under homotopy. The final equivalence follows by applying the first statement to
both inclusions.
\end{proof}
